\section*{3. The Beginning of a Friendship}

Toward the end of September 1909, Hermann Cohen spent a few days in Paris. For about a year, he had been corresponding with Meyerson, who, alongside his interest in epistemology, was involved in the administration of the Jewish Colonization Association (JCA).\footnote{Meyerson may already have been familiar with Cohen's work while composing \emph{Identité et réalité} (see, e.g., Meyerson 1908a, 209; tr. 1930a, 229). In March 1908, encouraged by the general director of the JCA, Sigismond Sonnenfeld, Meyerson sent Cohen a copy of \emph{Identité et réalité} in the hope of securing a review: ``Darf ich hinzufuegen, dass ich mit aengstlicher Spannung dem Urtheil des Altmeisters der deutschen Metaphysik ueber diesen Versuch entgegen sehe?'' (Meyerson to Cohen, March 13, 1908). This letter, like the rest of the correspondence between Cohen and Meyerson, is held at the Central Zionist Archives.} On September 30, Cohen sent Meyerson a postcard asking whether he would be available to meet the following afternoon. In the same postcard, Cohen specified that one of his former students, Ernst Cassirer, was also in Paris and was to take part in the meeting. The following day, October 1, 1909, at around 2 p.m., Cohen, Meyerson, and Cassirer met at the Café de la Régence.\footnote{See Cohen to Meyerson, September 30, 1909: ``Denken Sie nur: Dr. Cassirer ist hier. Können Sie morgen Donnerstag {[}\ldots{]} {[}um{]} 2 Uhr im Café Regence mit uns zusammen sein? {[}\ldots{]} Oder Freitag zur selben Zeit?'' For further evidence concerning Cohen's stay in Paris, see also Meyerson (2009, 345); Hartmann and Heimsoeth (1978, 45).}

Cassirer and Meyerson had not previously known each other, and their first encounter in Paris, mediated by Cohen, marked the beginning of a long-lasting friendship. After the meeting, Meyerson received a copy of Cassirer's \emph{Das Erkenntnisproblem in der Philosophie und Wissenschaft der neueren Zeit} (1906, 1907), a two-volume work devoted to the investigation of modern philosophy and science from the Renaissance through Kant. A letter to Cassirer from early November 1909 captures Meyerson's first reaction:

\begin{quote}
I have just received both volumes of your beautiful work and I hereby express to you my sincere thanks. I could not resist the temptation to immediately take a look inside and now I am doubly pleased by the enjoyment that the study of this significant book promises me. At the same time, nonetheless, I deeply regret not having known it before the publication of my own work. {[}\ldots{]} I remember with great pleasure the few hours we spent here together.\footnote{Meyerson to Cassirer, November 2, 1909 (ECN 18 DVD, doc. 121): ``Soeben erhielt ich die beiden Baende Ihres schoenen Werkes und spreche Ihnen hiermit meinen besten Dank aus. Ich habe der Versuchung nicht widerstehen koennen, sogleich einen Blick hineinzuwerfen und freue mich jetzt doppelt auf das Vergnuegen, das mir das Studium dieses bedeutenden Buches verspricht, indem ich zugleich tief bedaure dasselbe vor der Veroeffentlichung des meinigen nicht gekannt zu haben. {[}\ldots{]} Mit grossem Vergnuegen denke ich an die wenigen, hier gemeinsam verlebten Stunden zurueck.''}
\end{quote}

In the early twentieth century, both Cassirer and Meyerson were engaged in intense philosophical activity. Cassirer, then \emph{Privatdozent} at the University of Berlin, was occupied with revising the first edition of the \emph{Erkenntnisproblem} and writing \emph{Substanzbegriff und Funktionsbegriff} (1910) which was initially conceived as its third systematic volume. Meyerson, on the other hand, had recently published his first major work, \emph{Identité et réalité} (1908a), in which---in contrast to Cassirer---he defended a metaphysical reading of the foundations of science. Against this background, each offered valuable feedback on the other's research.

By late December 1909, Meyerson had already completed the ``first reading'' of the \emph{Erkenntnisproblem}, and he had contacted Xavier Léon---editor and co-founder of the \emph{Revue de métaphysique et de morale}---with the intention of reviewing Cassirer's work. Meyerson wrote to Cassirer:

\begin{quote}
I have now finished the study of your work---at least in its ``first reading,'' which will likely be followed by others. That I disagree with some passages, and precisely with some fundamental ones, will not surprise you, but this does not prevent me from feeling the greatest admiration for the whole work, for your astonishing knowledge, and for the extraordinary philosophical power with which you make use of it. For a more detailed judgment, however, you will have to wait until it appears in print. It has been arranged that my friend X. Léon will include a longer article (a sort of \emph{\underline{revue critique}}) on your book. I am not sure when it can appear, not even \underline{whether} it will actually appear at all.\footnote{Meyerson to Cassirer, December 23, 1909 (ECN 18 DVD, doc. 132): ``Mit dem Studium Ihres Werkes bin ich jetzt zu Ende---wenigstens `in erster Lesung', der wohl nach andere folgen werden. Dass ich mit Manchem, und zwar gerade mit manchem Grundsaetzlichen nicht einverstanden bin, wird Sie nicht ueberraschen, das hindert aber nicht, dass ich fuer das Ganze, fuer Ihr stupendes Wissen, sowie fuer die gewaltige philosophische Kraft mit der Sie dasselbe ausnutzen, die groesste Bewunderung empfinde. Auf ein naeheres Urtheil muessen Sie allerdings warten, bis es im Druck erscheint. Es ist naemlich mit Freund X. Léon abgemacht, dass er einen laengeren Artikel (eine sag. \underline{revue critique}) ueber das Buch aufnimmt. Wann derselbe erscheinen kann, bin ich aber nicht sicher, nicht einmal \underline{ob} er wirklich erscheinen wird.'' The underlining, both here and in the quotations that follow, is by Meyerson.}
\end{quote}

Meyerson's ``\emph{revue critique},'' which will be discussed more extensively in Section 5.2, was published in January 1911. From the outset, Cassirer welcomed Meyerson's proposal with great enthusiasm:

\begin{quote}
I am very pleased by what you said about my book. I know how deeply you yourself are involved in the problems with which my book is concerned, and your judgment is therefore of very special value to me. Regarding your intention to publish a review, I am, of course, especially grateful.\footnote{Cassirer to Meyerson, January 1, 1910 (ECN 18 DVD, doc. 134): ``Was Sie mir über mein Buch sagen, war mir eine große Freude; ich weiß, wie sehr Sie selbst mitten in den Problemen stehen, mit denen mein Buch sich beschäftigt und Ihr Urteil ist mir daher von ganz besonderem Wert. Was Ihre Absicht einer öffentlichen Besprechung betrifft, so bin ich Ihnen dafür natürlich ganz besonders dankbar.''}
\end{quote}

These words make it clear that Cassirer attached great importance to Meyerson's feedback. Since Meyerson's epistemological standpoint was radically opposed to Cassirer's, it represented an excellent test of the consistency of the \emph{Erkenntnisproblem}, and particularly of its systematic assumptions. Moreover, as stated in a letter to Meyerson from July 1910, Cassirer was keen to disseminate his ideas among members of the French philosophical community because, unlike its German counterpart, it was composed mainly of physicists and mathematicians.\footnote{See Cassirer to Meyerson, July 19, 1910 (ECN 18 DVD, doc. 150): ``die Verständigung mit dem französischen philosophischen Publikum, das ja mehr als bei uns in Deutschland ein Publikum von Mathematikern und Naturwissenschaftlern ist, liegt mir besonders am Herzen.''}

Even with Cassirer's support, Meyerson's work proceeded rather slowly. In April 1910, after completing the ``second reading'' of the \emph{Erkenntnisproblem}, Meyerson announced a visit to Berlin:

\begin{quote}
I hereby faithfully inform you that I intend to be in Berlin in about eight days. I am very much looking forward to a conversation with you---I have now also finished the ``second reading'' of your beautiful book and we will have a lot to discuss. I expect to arrive on Monday the 25th, around 9 a.m., and would like to know when I would least disturb you. I will probably depart on Tuesday evening at 11 p.m.\footnote{Meyerson to Cassirer, April 18, 1910 (ECN 18 DVD, doc. 141): ``Hiermit die ergebene Anzeige, dass ich in etwa acht Tagen in Berlin zu sein gedenke. Auf eine Unterhaltung mit Ihnen freue ich mich sehr---ich bin jetzt mit Ihrem schoenen Buche auch `in zweiter Lesung' fertig und wir werden vieles zu besprechen haben. Voraussichtlich komme ich Montag 25. vormittag von 9 an und moechte gern wissen, wann ich Sie am wenigsten stoere. Wahrscheinlich werde ich Dienstag abend um 11 abreisen.''}
\end{quote}

In the days between April 25 and May 3, 1909, Cassirer welcomed Meyerson into his home on more than one occasion. The correspondence does not reveal the content of their conversations. Nevertheless, it seems legitimate to suppose that they set out their respective viewpoints and anticipated some of the criticisms that would later appear in print. In particular, Meyerson might already have objected that \emph{Das Erkenntnisproblem} had completely neglected the role of chemistry in the historical development of the ``problem of knowledge'' in early modern science. Cassirer's exclusive focus on the history of physics had led him to overemphasize the role of function-concepts at the expense of the substance-concept, which figures much more prominently in the history of chemistry. However, by that time, Cassirer had probably already written the rather extensive section on chemistry that appears in Chapter 4 of \emph{Substanzbegriff und Funktionsbegriff---}a work initially conceived as a systematic counterpart to \emph{Das Erkenntnisproblem.} Far from neglecting the history of chemistry, Cassirer had tried to show that it was dominated by the very same model of concept formation (\emph{Begriffsbildung}) that he had derived from his analysis of the history of physics. A few remarks, possibly added at a later stage in the drafting of \emph{Substanzbegriff und Funktionsbegriff}, suggest that Cassirer regarded Meyerson's reliance on the substance-concept not as a consequence of his focus on the history of chemistry but rather as the result of a mistaken conception of scientific concept formation that shaped his historiographical work.
